% GENERATED FILE. Edit canonical lessons, then run npm run generate:books.
% canonical-source-hash: fnv1a64:08c1e4d93084de17
% canonical-lessons: MR-C164-jahir-karne, MR-C164-talya-vajavne, MR-C164-olandne, MR-C164-vadhavne, MR-C164-kami-karne

\chapter{To announce, To clap, To cross, To increase, To reduce}
\label{ch:mr-to-announce-to-clap-to-cross-to-increase-to-reduce}
\hlchaptermodality{164}

\begin{chapteropening}
\textbf{By the end of this chapter:} \emph{I can say to announce, to clap, to cross, to increase and to reduce.}

Complete the last lesson of chapter 164: I can say to announce, to clap, to cross, to increase and to reduce.
\end{chapteropening}

\section[\texorpdfstring{\mr{जाहीर} \mr{करणे} (jāhīr karṇe)}{जाहीर करणे (jāhīr karṇe)}]{\mr{जाहीर} \mr{करणे} (jāhīr karṇe) — to announce}

\label{lesson:MR-C164-jahir-karne}



\begin{quote}
Before the new one: say the Marathi for to check, then the Marathi for to sign.
\end{quote}

\subsection*{You'll want to know: \mr{जाहीर} \mr{करणे}}
\textbf{\mr{जाहीर} \mr{करणे}} — \emph{jāhīr karṇe} — \textquotedblleft{}to announce\textquotedblright{}.

\begin{grammarlens}[title={What you've built}]
One more word for this chapter.
\end{grammarlens}

\subsection*{Your turn}
\begin{itemize}
\raggedright
  \item \emph{Say it:} \emph{jāhīr karṇe}
  \item \emph{Say it:} \emph{jāhīr karṇe}, once more
  \item \emph{Say it:} say \emph{jāhīr karṇe}
  \item \emph{Recall:} say the Marathi for to check, then the Marathi for to sign, then say \emph{jāhīr karṇe} again
\end{itemize}

\begin{tcolorbox}[breakable,colback=teal!4,colframe=teal!35!black,title={Before you move on}]
What does \mr{जाहीर} \mr{करणे} mean? (\textquotedblleft{}To announce\textquotedblright{}.) Say it once more.
\end{tcolorbox}
\section[\texorpdfstring{\mr{टाळ्या} \mr{वाजवणे} (ṭāḷyā vājavṇe)}{टाळ्या वाजवणे (ṭāḷyā vājavṇe)}]{\mr{टाळ्या} \mr{वाजवणे} (ṭāḷyā vājavṇe) — to clap}

\label{lesson:MR-C164-talya-vajavne}



\begin{quote}
Before the new one: say the Marathi for to sign, then the Marathi for to announce.
\end{quote}

\subsection*{You'll want to know: \mr{टाळ्या} \mr{वाजवणे}}
\textbf{\mr{टाळ्या} \mr{वाजवणे}} — \emph{ṭāḷyā vājavṇe} — \textquotedblleft{}to clap\textquotedblright{}.

\begin{grammarlens}[title={What you've built}]
One more word for this chapter.
\end{grammarlens}

\subsection*{Your turn}
\begin{itemize}
\raggedright
  \item \emph{Say it:} \emph{ṭāḷyā vājavṇe}
  \item \emph{Say it:} \emph{ṭāḷyā vājavṇe}, once more
  \item \emph{Say it:} say \emph{ṭāḷyā vājavṇe}
  \item \emph{Recall:} say the Marathi for to sign, then the Marathi for to announce, then say \emph{ṭāḷyā vājavṇe} again
\end{itemize}

\begin{tcolorbox}[breakable,colback=teal!4,colframe=teal!35!black,title={Before you move on}]
What does \mr{टाळ्या} \mr{वाजवणे} mean? (\textquotedblleft{}To clap\textquotedblright{}.) Say it once more.
\end{tcolorbox}
\section[\texorpdfstring{\mr{ओलांडणे} (olāṇḍṇe)}{ओलांडणे (olāṇḍṇe)}]{\mr{ओलांडणे} (olāṇḍṇe) — to cross}

\label{lesson:MR-C164-olandne}



\begin{quote}
Before the new one: say the Marathi for to announce, then the Marathi for to clap.
\end{quote}

\subsection*{You'll want to know: \mr{ओलांडणे}}
\textbf{\mr{ओलांडणे}} — \emph{olāṇḍṇe} — \textquotedblleft{}to cross\textquotedblright{}.

\begin{grammarlens}[title={What you've built}]
One more word for this chapter.
\end{grammarlens}

\subsection*{Your turn}
\begin{itemize}
\raggedright
  \item \emph{Say it:} \emph{olāṇḍṇe}
  \item \emph{Say it:} \emph{olāṇḍṇe}, once more
  \item \emph{Say it:} say \emph{olāṇḍṇe}
  \item \emph{Recall:} say the Marathi for to announce, then the Marathi for to clap, then say \emph{olāṇḍṇe} again
\end{itemize}

\begin{tcolorbox}[breakable,colback=teal!4,colframe=teal!35!black,title={Before you move on}]
What does \mr{ओलांडणे} mean? (\textquotedblleft{}To cross\textquotedblright{}.) Say it once more.
\end{tcolorbox}
\section[\texorpdfstring{\mr{वाढवणे} (vāḍhavṇe)}{वाढवणे (vāḍhavṇe)}]{\mr{वाढवणे} (vāḍhavṇe) — to increase}

\label{lesson:MR-C164-vadhavne}



\begin{quote}
Before the new one: say the Marathi for to clap, then the Marathi for to cross.
\end{quote}

\subsection*{You'll want to know: \mr{वाढवणे}}
\textbf{\mr{वाढवणे}} — \emph{vāḍhavṇe} — \textquotedblleft{}to increase\textquotedblright{}.

\begin{grammarlens}[title={What you've built}]
One more word for this chapter.
\end{grammarlens}

\subsection*{Your turn}
\begin{itemize}
\raggedright
  \item \emph{Say it:} \emph{vāḍhavṇe}
  \item \emph{Say it:} \emph{vāḍhavṇe}, once more
  \item \emph{Say it:} say \emph{vāḍhavṇe}
  \item \emph{Recall:} say the Marathi for to clap, then the Marathi for to cross, then say \emph{vāḍhavṇe} again
\end{itemize}

\begin{tcolorbox}[breakable,colback=teal!4,colframe=teal!35!black,title={Before you move on}]
What does \mr{वाढवणे} mean? (\textquotedblleft{}To increase\textquotedblright{}.) Say it once more.
\end{tcolorbox}
\section[\texorpdfstring{\mr{कमी} \mr{करणे} (kamī karṇe)}{कमी करणे (kamī karṇe)}]{\mr{कमी} \mr{करणे} (kamī karṇe) — to reduce}

\label{lesson:MR-C164-kami-karne}



\begin{quote}
Before the new one: say the Marathi for to cross, then the Marathi for to increase.
\end{quote}

\subsection*{You'll want to know: \mr{कमी} \mr{करणे}}
\textbf{\mr{कमी} \mr{करणे}} — \emph{kamī karṇe} — \textquotedblleft{}to reduce\textquotedblright{}.

\begin{grammarlens}[title={What you've built}]
One more word for this chapter.
\end{grammarlens}

\subsection*{Your turn}
\begin{itemize}
\raggedright
  \item \emph{Say it:} \emph{kamī karṇe}
  \item \emph{Say it:} \emph{kamī karṇe}, once more
  \item \emph{Say it:} say \emph{kamī karṇe}
  \item \emph{Recall:} say the Marathi for to cross, then the Marathi for to increase, then say \emph{kamī karṇe} again
\end{itemize}

\begin{tcolorbox}[breakable,colback=teal!4,colframe=teal!35!black,title={Before you move on}]
What does \mr{कमी} \mr{करणे} mean? (\textquotedblleft{}To reduce\textquotedblright{}.) Say it once more.
\end{tcolorbox}
